\section{Algebra structure}

% Source-numbering note: arXiv:1606.00608 uses one theorem counter for
% definitions, propositions, theorems, examples, and observations in Section 4.
% The three examples after Theorem 4.9 are numbered 4.10--4.12; hence
% `thm:IV.13` is Theorem 4.14 in the arXiv source.

\begin{definition}[Algebra-structure data for an MPDO]%
    \label{def:mpdo_algebra_structure_data}
    \lean{MPOTensor.AlgebraStructureData}
    \leanok
    \uses{def:mpo_tensor}
    A family of support $*$-algebras $\{\mathcal{A}_n\}_{n \ge 0}$, one at
    every blocking size, equipped with blocking maps
    $m_n : \mathcal{A}_n \otimes \mathcal{A}_n \to \mathcal{A}_{2n}$ and
    inclusion maps $\iota_n : \mathcal{A}_n \to \mathcal{A}_{n+1}$ such that,
    after viewing all elements inside the ambient bond-space matrix algebra,
    $m_n$ is ordinary matrix multiplication and $\iota_n$ is the ambient
    inclusion.
\end{definition}

\begin{definition}[Algebra-structure RFP predicate]%
    \label{def:mpdo_rfp_via_algebra}
    \lean{MPOTensor.IsRFP_MPDO_via_algebra}
    \leanok
    \uses{def:mpdo_algebra_structure_data, def:mpo_tensor}
    An MPO tensor $M$ satisfies the \emph{algebra-structure formulation} used
    here when there exist algebra-structure data whose support algebra at every
    positive blocking size coincides with the fixed-point algebra of the adjoint
    blocked transfer map. This is a nontrivial algebraic condition, but it is still
    weaker than the full coefficient formulation of
    \cite[Theorem~4.14(ii)]{Cirac2016MPDO_arXiv}; that formulation also
    includes the coefficient family $c_{\alpha,\beta,\gamma}^{(L)}$.
\end{definition}

\begin{theorem}[RFP yields a stationary algebra tower under a faithful fixed point]%
    \label{thm:mpdo_rfp_via_algebra_from_rfp}
    \lean{MPOTensor.isRFP_MPDO_via_algebra_of_isRFP_of_isTP_of_posDef_fixed}
    \leanok
    \uses{def:is_rfp_mpdo, def:mpdo_rfp_via_algebra}
    Assume the doubled tensor is trace-preserving and the transfer map admits a
    positive-definite fixed point. If the transfer map is idempotent, then there
    exists a stationary family of support algebras, all equal to the fixed-point
    algebra of the adjoint transfer map, with multiplication given by matrix
    multiplication and inclusions given by the identity.
\end{theorem}
\begin{proof}\leanok
    Under the trace-preserving normalization and the faithful fixed point, Wolf's
    fixed-point-algebra theorem gives a $*$-subalgebra of adjoint fixed points.
    Idempotence identifies every positive blocked transfer map with the original
    one, so the same fixed-point algebra works at every blocking size.
\end{proof}

\begin{theorem}[Fusion yields a stationary algebra tower under a faithful fixed point]%
    \label{thm:mpdo_rfp_via_algebra_from_fusion}
    \lean{MPOTensor.isRFP_MPDO_via_algebra_of_isRFP_MPDO_via_fusion_of_isTP_of_posDef_fixed}
    \leanok
    \uses{thm:mpdo_rfp_of_fusion, thm:mpdo_rfp_via_algebra_from_rfp}
    Assume the doubled tensor is trace-preserving and the transfer map admits a
    positive-definite fixed point. If the transfer-map fusion criterion holds,
    then the MPO satisfies the algebra-structure formulation of the
    renormalization fixed-point condition.
\end{theorem}
\begin{proof}\leanok
    Combine Theorem~\ref{thm:mpdo_rfp_of_fusion} with
    Theorem~\ref{thm:mpdo_rfp_via_algebra_from_rfp}.
\end{proof}

\begin{theorem}[Stabilized blocked adjoint fixed points yield algebra data]%
    \label{thm:mpdo_algebra_from_stabilized_adjoint_fixedpoints}
    \lean{MPOTensor.isRFP_MPDO_via_algebra_of_adjointFixedPoints_eq_of_isTP_of_posDef_fixed}
    \leanok
    \uses{def:mpdo_rfp_via_algebra}
    Assume the doubled tensor is trace-preserving and the transfer map admits a
    positive-definite fixed point. If, for every positive blocked size $n$, the
    adjoint fixed-point space of the blocked transfer map agrees with the
    adjoint fixed-point space of the original transfer map, then the
    algebra-structure formulation holds.
\end{theorem}
\begin{proof}\leanok
    Use the adjoint fixed-point algebra of the original transfer map as a constant
    support-algebra tower. The stabilization hypothesis identifies that same
    algebra with the adjoint fixed-point space of every positive blocked transfer
    map, so the stationary family satisfies the definition.
\end{proof}

This does not give the full converse direction of
\cite[Theorem~4.14]{Cirac2016MPDO_arXiv}: the compatibility condition above
only sees stabilized adjoint fixed-point algebras, not the coefficient/BNT data.
In particular, \cite[Theorem~4.14(ii)]{Cirac2016MPDO_arXiv} also requires the special diagonal matrices
$\chi_{\alpha,\beta,\gamma}$. The coefficient data attached to the algebra tower are:
a chosen basis of each support algebra, coordinate maps into a finite scalar
family, and the corresponding multiplication / inclusion coefficient arrays.

\begin{definition}[Blocked coefficients for the algebra tower]%
    \label{def:mpdo_blocked_coefficients}
    \lean{MPOTensor.AlgebraStructureData.toBlockedCoefficients}
    \leanok
    \uses{def:mpdo_algebra_structure_data}
    For each blocked size $n$, once a basis of $\mathcal{A}_n$ is chosen,
    every element of $\mathcal{A}_n$ is identified with a finite coefficient
    family on that basis.
\end{definition}

\begin{theorem}[Reconstruction from blocked coefficients]%
    \label{thm:mpdo_reconstruct_blocked_coefficients}
    \lean{MPOTensor.AlgebraStructureData.reconstructFromBlockedCoefficients_apply}
    \leanok
    \uses{def:mpdo_blocked_coefficients}
    Reconstructing a coefficient family gives the corresponding basis
    expansion $\sum_i a_i b_i$ in $\mathcal{A}_n$.
\end{theorem}
\begin{proof}\leanok
    This is exactly the finite-basis reconstruction formula for the chosen basis of
    $\mathcal{A}_n$.
\end{proof}

\begin{definition}[Blocked multiplication coefficients]%
    \label{def:mpdo_blocked_structure_coefficients}
    \lean{MPOTensor.AlgebraStructureData.blockedStructureCoefficients}
    \leanok
    \uses{def:mpdo_blocked_coefficients}
    The blocking multiplication map $m_n$ determines a coefficient family for
    the product of any two chosen basis elements of $\mathcal{A}_n$, now viewed
    in the basis of $\mathcal{A}_{2n}$.
\end{definition}

\begin{theorem}[Blocked multiplication reconstructs from its coefficients]%
    \label{thm:mpdo_mul_coefficients_reconstruct}
    \lean{MPOTensor.AlgebraStructureData.coe_mul_eq_sum_blockedStructureCoefficients}
    \leanok
    \uses{def:mpdo_blocked_structure_coefficients}
    In the ambient bond-space matrix algebra, the product of two chosen basis
    elements of $\mathcal{A}_n$ is the sum of the basis elements of
    $\mathcal{A}_{2n}$ weighted by these extracted coefficients.
\end{theorem}
\begin{proof}\leanok
    Apply the reconstruction formula of
    Theorem~\ref{thm:mpdo_reconstruct_blocked_coefficients} to the element
    $m_n(b_i,b_j) \in \mathcal{A}_{2n}$.
\end{proof}

\begin{theorem}[Compatible adjoint fixed points reconstruct from blocked coefficients]%
    \label{thm:mpdo_fixedpoint_coefficients_reconstruct}
    \lean{MPOTensor.AlgebraStructureData.reconstructFromBlockedCoefficients_of_fixedPoint}
    \leanok
    \uses{def:mpdo_rfp_via_algebra, def:mpdo_blocked_coefficients}
    If the algebra tower is compatible with an MPO tensor $M$, then every adjoint
    blocked fixed point of the blocked transfer map $E_n$ belongs to
    $\mathcal{A}_n$ and is recovered exactly from its extracted coefficient
    family.
\end{theorem}
\begin{proof}\leanok
    Compatibility identifies $\mathcal{A}_n$ with the adjoint fixed-point algebra of
    $E_n$, so the fixed point first becomes an element of $\mathcal{A}_n$ and then
    Theorem~\ref{thm:mpdo_reconstruct_blocked_coefficients} reconstructs it.
\end{proof}

\begin{theorem}[Adjoint fixed points descend through the algebra tower]%
    \label{thm:mpdo_adjoint_fix_descent}
    \lean{MPOTensor.adjoint_transferMap_apply_of_isRFP_MPDO_via_algebra}
    \leanok
    \uses{def:mpdo_rfp_via_algebra}
    If an MPO tensor satisfies the algebra-structure formulation, then every
    adjoint fixed point of the blocked transfer map $E_n$ at a positive
    blocking size $n$ is already an adjoint fixed point of the unblocked
    transfer map $E_1$.
\end{theorem}
\begin{proof}\leanok
    Compatibility at size $n$ places $X$ in $\mathcal{A}_n$. The inclusion map
    $\iota_n$ lifts $X$ into $\mathcal{A}_{n+1}$, and since $\iota_n$ is just an
    inclusion it does not change the underlying operator, so we continue to
    denote the lifted element by $X$. Thus compatibility at size $n + 1$ yields
    $E_{n+1}^{\dagger} X = X$. Rewriting $E_{n+1} = E_n \circ E_1$ and applying
    $(A \circ B)^{\dagger} = B^{\dagger} \circ A^{\dagger}$ together with the
    hypothesis $E_n^{\dagger} X = X$ extracts $E_1^{\dagger} X = X$.
\end{proof}

\begin{lemma}[Reverse inclusion of adjoint fixed points]%
    \label{lem:mpdo_adjoint_fix_reverse}
    \lean{MPOTensor.adjoint_blockedTransferMap_apply_of_adjoint_transferMap_apply}
    \leanok
    Every adjoint fixed point of the unblocked transfer map $E_1$ is an
    adjoint fixed point of every blocked transfer map $E_n$.
\end{lemma}
\begin{proof}\leanok
    Induction on $n$ using $E_{n+1} = E_n \circ E_1$ and
    $(A \circ B)^{\dagger} = B^{\dagger} \circ A^{\dagger}$; the base case is
    $E_0 = \mathrm{id}$. This step does not require the algebra-structure data.
\end{proof}

\begin{lemma}[Blocked powers of adjoint eigenvectors]%
    \label{lem:mpdo_adjoint_blocked_eigenvector}
    \lean{MPOTensor.adjoint_blockedTransferMap_apply_of_adjoint_transferMap_eigenvector}
    \leanok
    If $E=E_1$ and $E^\dagger X=\lambda X$, then
    $E_n^\dagger X=\lambda^n X$ for every $n$.
\end{lemma}
\begin{proof}\leanok
    The case $n=0$ is $E_0=\mathrm{id}$. For the induction step,
    $E_{n+1}=E_n\circ E_1$, hence
    \[
        E_{n+1}^{\dagger}X
        = E_1^{\dagger}E_n^{\dagger}X
        = E_1^{\dagger}(\lambda^n X)
        = \lambda^n E_1^{\dagger}X
        = \lambda^{n+1}X.
    \]
\end{proof}

\begin{theorem}[Finite-order adjoint eigenvalues for algebra towers]%
    \label{thm:mpdo_algebra_root_of_unity_eigenvalue}
    \lean{MPOTensor.adjoint_transferMap_eigenvalue_eq_one_of_isRFP_MPDO_via_algebra}
    \leanok
    \uses{def:mpdo_rfp_via_algebra, thm:mpdo_adjoint_fix_descent,
        lem:mpdo_adjoint_blocked_eigenvector}
    Assume the algebra-structure formulation holds. If $X\ne 0$,
    $E_1^\dagger X=\lambda X$, and $\lambda^n=1$ for some $n>0$, then
    $\lambda=1$.
\end{theorem}
\begin{proof}\leanok
    By Lemma~\ref{lem:mpdo_adjoint_blocked_eigenvector},
    $E_n^\dagger X=\lambda^n X=X$. Theorem~\ref{thm:mpdo_adjoint_fix_descent} gives
    $E_1^\dagger X=X$. Thus $\lambda X=X$, so $(\lambda-1)X=0$. Since
    $X\ne0$, one obtains $\lambda=1$.
\end{proof}

\begin{theorem}[Adjoint fixed-point equality for algebra towers]%
    \label{thm:mpdo_adjoint_fix_equality_from_algebra}
    \lean{MPOTensor.adjoint_blockedTransferMap_apply_iff_of_isRFP_MPDO_via_algebra}
    \leanok
    \uses{def:mpdo_rfp_via_algebra, thm:mpdo_adjoint_fix_descent,
        lem:mpdo_adjoint_fix_reverse}
    If an MPO tensor satisfies the algebra-structure formulation, then for
    every positive blocking size $n$ the adjoint fixed points of $E_n$ are
    exactly the adjoint fixed points of $E_1$.
\end{theorem}
\begin{proof}\leanok
    Combine Theorem~\ref{thm:mpdo_adjoint_fix_descent} with
    Lemma~\ref{lem:mpdo_adjoint_fix_reverse}.
\end{proof}

\begin{theorem}[Algebra data give the stationary fixed-point tower]%
    \label{thm:mpdo_stationary_tower_from_algebra}
    \lean{MPOTensor.stationaryOfFaithfulFixedPoint_compatible_of_isRFP_MPDO_via_algebra}
    \leanok
    \uses{def:mpdo_rfp_via_algebra,
        thm:mpdo_adjoint_fix_equality_from_algebra}
    Under the trace-preserving normalization and a positive-definite fixed
    point, any algebra-structure witness forces the stationary adjoint
    fixed-point tower itself to be compatible with the MPO tensor.
\end{theorem}
\begin{proof}\leanok
    Apply the adjoint fixed-point equality obtained from the algebra-structure
    predicate to the stationary tower constructed from the faithful fixed point.
\end{proof}

\begin{remark}[Why algebra data do not imply fusion data]%
    \label{rem:mpdo_algebra_to_fusion_gap}
    \uses{def:mpdo_rfp_via_algebra, def:mpdo_rfp_via_fusion, thm:newton_girard}
    Write $E=E_1$ for the one-site transfer map. The present algebra predicate
    gives, for every $n>0$,
    \[
        \operatorname{Fix}((E^n)^\dagger)
        = \operatorname{Fix}(E^\dagger).
    \]
    Hence a vector $X$ with $E^\dagger X=\lambda X$ and $\lambda^n=1$ lies in
    $\operatorname{Fix}(E^\dagger)$, so $(\lambda-1)X=0$ and $\lambda=1$ by
    Theorem~\ref{thm:mpdo_algebra_root_of_unity_eigenvalue}. It does not imply
    $E^2=E$. For instance, if $\Pi_{\mathrm{diag}}$ is the projection onto
    diagonal matrices and $0<\varepsilon<1$, set
    \[
        E(X)
        = \varepsilon\, X
        + (1 - \varepsilon)\, \Pi_{\mathrm{diag}}(X),
    \]
    so that
    \[
        E^n(X)
        = \Pi_{\mathrm{diag}}(X)
        + \varepsilon^n(X-\Pi_{\mathrm{diag}}(X)),
        \qquad
        E^2-E
        = (\varepsilon^2-\varepsilon)
        (\mathrm{id}-\Pi_{\mathrm{diag}}).
    \]
    The fixed-point algebra is diagonal for every $n>0$, but $E^2\ne E$.

    The missing source step is the coefficient argument in
    \cite[Appendix~C.4]{Cirac2016MPDO_arXiv}. One needs the same-length
    product law
    \begin{equation}\label{eq:mpdo_chi_product_law}
        O_L(M_\alpha)O_L(M_\beta)
        =
        \sum_\gamma
        \tr(\chi_{\alpha,\beta,\gamma}^{\,L})
        O_L(M_\gamma)
    \end{equation}
    and the length-one idempotent trace identity
    \begin{equation}\label{eq:mpdo_idempotent_trace_identity}
        m_\gamma
        =
        \sum_{\alpha,\beta}
        \tr(\chi_{\alpha,\beta,\gamma})
        m_\alpha m_\beta.
    \end{equation}
    The proof of \cite[Theorem~4.14]{Cirac2016MPDO_arXiv} uses
    (\ref{eq:mpdo_chi_product_law}) and (\ref{eq:mpdo_idempotent_trace_identity}),
    positivity, and the simple-matrix lemma
    \cite[lines~1155--1163]{Cirac2016MPDO_arXiv} to give
    \begin{equation}\label{eq:mpdo_nu_chi_witness}
        \{\nu_{\gamma,k}\}_k
        =
        \{\mu_{\alpha,k_1}\mu_{\beta,k_2}
        \chi_{\alpha,\beta,\gamma,k_3}\}_{\alpha,\beta,k_1,k_2,k_3},
        \qquad
        \tr(\nu_\gamma)=m_\gamma.
    \end{equation}
    The reconstruction maps in \cite[lines~2065--2085]{Cirac2016MPDO_arXiv}
    then use
    \begin{equation}\label{eq:mpdo_reconstruction_maps}
        \widetilde R_2\left(\bigoplus_\gamma X_\gamma\right)
        = \bigoplus_\gamma \frac{\nu_\gamma}{m_\gamma}\otimes X_\gamma,
        \qquad
        \widetilde R_1\left(\bigoplus_\gamma X_\gamma\right)
        = \bigoplus_\gamma \frac{\mu_\gamma}{m_\gamma}\otimes X_\gamma.
    \end{equation}
    Thus $\tr(\nu_\gamma)=m_\gamma$ and the defining equality
    $m_\gamma=\tr(\mu_\gamma)$ normalize the factors
    $\nu_\gamma/m_\gamma$ and $\mu_\gamma/m_\gamma$, respectively, in
    $T=\widetilde R_2\widetilde T R_1$ and
    $S=\widetilde R_1\widetilde S R_2$. The scalar
    Newton--Girard power-sum identity needed to obtain the trace-power form is
    already available as
    Theorem~\ref{thm:newton_girard}; what remains is the MPDO construction of
    the BNT-label witness carrying (\ref{eq:mpdo_nu_chi_witness}) and
    (\ref{eq:mpdo_reconstruction_maps}).
\end{remark}

\subsection{Diagonal \texorpdfstring{$\chi$}{chi}-matrices and the trace-power formula}

The following statements isolate the form of
\cite[Theorem~4.14(ii)]{Cirac2016MPDO_arXiv}: the structure coefficients
of the length-$L$ operator algebra are trace-powers of positive diagonal
matrices $\chi_{\alpha,\beta,\gamma}$. They state the coefficient identity and
the elementary trace-power identity for diagonal matrices.

\begin{definition}[Diagonal $\chi$ family]%
    \label{def:mpdo_diagonal_chi_family}
    \lean{MPOTensor.DiagonalChiFamily}
    \leanok
    A family of diagonal complex matrices $\chi_{\alpha,\beta,\gamma}$ indexed
    by ordered triples drawn from a common index type, each of a specified
    size. This is the matrix $\chi_{\alpha,\beta,\gamma}$ of
    \cite[Theorem~4.14(ii)]{Cirac2016MPDO_arXiv}.
\end{definition}

\begin{theorem}[Trace of $\chi^L$ equals the sum of $L$-th powers of entries]%
    \label{thm:mpdo_trace_matrix_pow}
    \lean{MPOTensor.DiagonalChiFamily.trace_matrix_pow}
    \leanok
    \uses{def:mpdo_diagonal_chi_family}
    For every triple $(\alpha, \beta, \gamma)$, if $\chi_{\alpha,\beta,\gamma}$
    has size $r_{\alpha,\beta,\gamma}$, then for every $L \ge 0$,
    \[
        \operatorname{tr}(\chi_{\alpha,\beta,\gamma}^{L})
        = \sum_{k=1}^{r_{\alpha,\beta,\gamma}}
        \chi_{\alpha,\beta,\gamma,k}^{L},
    \]
    where $\chi_{\alpha,\beta,\gamma,k}$ are the diagonal entries.
\end{theorem}
\begin{proof}\leanok
    Diagonal matrices raise to powers entrywise and the trace of a diagonal matrix
    is the sum of its entries.
\end{proof}

\begin{definition}[Trace-power form of a structure-coefficient family]%
    \label{def:mpdo_has_chi_trace_power_form}
    \lean{MPOTensor.HasChiTracePowerForm}
    \leanok
    \uses{def:mpdo_diagonal_chi_family}
    A binary compatibility predicate between an abstract structure-coefficient
    family $c^{(L)}_{\alpha,\beta,\gamma}$ and a diagonal $\chi$ family: the
    pair $(c, \chi)$ is said to be in \emph{trace-power form} when
    \[
        c^{(L)}_{\alpha,\beta,\gamma}
        = \operatorname{tr}(\chi_{\alpha,\beta,\gamma}^{L})
        = \sum_{k=1}^{r_{\alpha,\beta,\gamma}}
        \chi_{\alpha,\beta,\gamma,k}^{L}
    \]
    for every $L \ge 0$ and every triple $(\alpha, \beta, \gamma)$, where
    $r_{\alpha,\beta,\gamma}$ is the size of $\chi_{\alpha,\beta,\gamma}$. This
    is the formal analog of the target identity of
    \cite[Theorem~4.14(ii)]{Cirac2016MPDO_arXiv}; the existential
    version---``there exists $\chi$ for which $(c,\chi)$ has trace-power
    form''---is obtained by quantifying over $\chi$.
\end{definition}

\begin{definition}[BNT-label coefficient family]%
    \label{def:mpdo_bnt_label_coefficient_family}
    \lean{MPOTensor.BNTLabelCoefficientFamily}
    \lean{MPOTensor.BNTLabelCoefficientFamily.ofChi}
    \lean{MPOTensor.BNTLabelCoefficientFamily.ofChi_coeff}
    \lean{MPOTensor.BNTLabelCoefficientFamily.ofChi_coeff_eq_trace_matrix_pow}
    \leanok
    The same-length coefficient system
    $c^{(L)}_{\alpha,\beta,\gamma}$ from \cite[Theorem~4.14(ii)]{Cirac2016MPDO_arXiv}, indexed by
    fixed BNT labels $\alpha,\beta,\gamma$. The coefficient is attached to
    the length-$L$ product
    $O_L(M_\alpha)O_L(M_\beta)$, not to the blocked-basis multiplication
    $\mathcal A_n\times\mathcal A_n\to\mathcal A_{2n}$.
    When a diagonal $\chi$-family has already been constructed, it also
    canonically determines the coefficient family
    $c^{(L)}_{\alpha,\beta,\gamma}
        = \operatorname{tr}(\chi_{\alpha,\beta,\gamma}^{L})
        = \sum_r\chi_{\alpha,\beta,\gamma,r}^{\,L}$, as in
    \cite[Appendix~C.4]{Cirac2016MPDO_arXiv}.
\end{definition}

\begin{definition}[BNT-label operator family]%
    \label{def:mpdo_bnt_label_operator_family}
    \lean{MPOTensor.BNTLabelOperatorFamily}
    \leanok
    A BNT-label operator family records, for each chain length $L$, the
    operators $O_L(M_\alpha)$ indexed by the fixed BNT labels $\alpha$.
    The ambient algebra at length $L$ is left abstract here; the later
    comparison step must relate it to the chosen blocked support algebras.
\end{definition}

\begin{definition}[Same-length BNT product form]%
    \label{def:mpdo_bnt_label_same_length_product_form}
    \lean{MPOTensor.BNTLabelOperatorFamily.HasSameLengthProductForm}
    \leanok
    \uses{def:mpdo_bnt_label_operator_family,
        def:mpdo_bnt_label_coefficient_family}
    A BNT-label operator family and a BNT-label coefficient system have
    same-length product form when, for every positive length $L$,
    \[
        O_L(M_\alpha)O_L(M_\beta)
        = \sum_\gamma
        c^{(L)}_{\alpha,\beta,\gamma} O_L(M_\gamma).
    \]
    This is the product identity appearing in \cite[Theorem~4.14(ii)]{Cirac2016MPDO_arXiv}, stated
    independently of the blocked-basis multiplication
    $\mathcal A_n\times\mathcal A_n\to\mathcal A_{2n}$.
\end{definition}

\begin{theorem}[Same-length BNT product expansion]%
    \label{thm:mpdo_bnt_label_same_length_product_eq_sum}
    \lean{MPOTensor.BNTLabelOperatorFamily.HasSameLengthProductForm.eq_sum}
    \leanok
    \uses{def:mpdo_bnt_label_same_length_product_form}
    Under same-length BNT product form, the product of two length-$L$
    BNT-label operators is the corresponding finite linear combination of
    length-$L$ BNT-label operators.
\end{theorem}
\begin{proof}\leanok
    This is exactly the defining equality of same-length product form.
\end{proof}

\begin{definition}[BNT-label trace-scalar family]%
    \label{def:mpdo_bnt_label_trace_scalar_family}
    \lean{MPOTensor.BNTLabelTraceScalarFamily}
    \leanok
    A BNT-label trace-scalar family records the scalars
    $m_\alpha=\operatorname{tr}(\mu_\alpha)$ indexed by the fixed BNT labels
    $\alpha$, as in the idempotent condition of \cite[Theorem~4.14(ii)]{Cirac2016MPDO_arXiv}.
\end{definition}

\begin{definition}[BNT idempotent coefficient form]%
    \label{def:mpdo_bnt_label_idempotent_coefficient_form}
    \lean{MPOTensor.BNTLabelTraceScalarFamily.HasIdempotentCoefficientForm}
    \leanok
    \uses{def:mpdo_bnt_label_trace_scalar_family,
        def:mpdo_bnt_label_coefficient_family}
    A BNT-label trace-scalar family and a BNT-label coefficient system have
    idempotent coefficient form when, for every BNT label $\gamma$,
    \[
        m_\gamma =
        \sum_{\alpha,\beta}
        c^{(1)}_{\alpha,\beta,\gamma}m_\alpha m_\beta.
    \]
    This is the idempotent condition in \cite[Theorem~4.14(ii)]{Cirac2016MPDO_arXiv}, stated separately
    from the trace-power formula.
\end{definition}

\begin{theorem}[BNT idempotent coefficient expansion]%
    \label{thm:mpdo_bnt_label_idempotent_coefficient_eq_sum}
    \lean{MPOTensor.BNTLabelTraceScalarFamily.HasIdempotentCoefficientForm.eq_sum}
    \leanok
    \uses{def:mpdo_bnt_label_idempotent_coefficient_form}
    Under BNT idempotent coefficient form, each $m_\gamma$ is the
    corresponding finite linear combination of products $m_\alpha m_\beta$
    with the length-one coefficients $c^{(1)}_{\alpha,\beta,\gamma}$.
\end{theorem}
\begin{proof}\leanok
    This is exactly the defining equality of BNT idempotent coefficient form.
\end{proof}

\begin{definition}[Blocked-basis BNT-label assignment]%
    \label{def:mpdo_bnt_blocked_basis_label_assignment}
    \lean{MPOTensor.BNTBlockedBasisLabelAssignment}
    \lean{MPOTensor.BNTBlockedBasisLabelAssignment.blockedLabel}
    \leanok
    \uses{def:mpdo_blocked_structure_coefficients}
    The BNT-label assignment consists, for every positive blocked length $n$,
    of a source label map $\sigma_n$ from the chosen basis of
    $\mathcal A_n$ and a target label map $\tau_n$ from the chosen basis of
    $\mathcal A_{2n}$ to the fixed BNT labels. It also gives the combined
    label map $\sigma_n \sqcup \tau_n$ on the disjoint union of these two
    blocked bases.
\end{definition}

\begin{definition}[BNT comparison with blocked coefficients]%
    \label{def:mpdo_bnt_blocked_basis_coefficient_comparison}
    \lean{MPOTensor.BNTBlockedBasisCoefficientComparison}
    \lean{MPOTensor.BNTBlockedBasisCoefficientComparison.sourceLabel}
    \lean{MPOTensor.BNTBlockedBasisCoefficientComparison.targetLabel}
    \leanok
    \uses{def:mpdo_bnt_label_coefficient_family,
        def:mpdo_blocked_structure_coefficients,
        def:mpdo_bnt_blocked_basis_label_assignment}
    A comparison between the BNT-label coefficients and the chosen
    blocked-basis coefficients consists of a BNT-label assignment together
    with the equality
    \[
        c^{(n)}_{i,j,k}
        =
        c^{(n)}_{\sigma_n(i),\sigma_n(j),\tau_n(k)}.
    \]
    This is a blocked-basis comparison: the left hand side is expanded in the
    chosen basis of $\mathcal A_{2n}$, so the statement is not itself the
    same-length BNT product law.
    This records the coefficient-comparison statement separately from the
    construction of the label maps from the MPDO tensor.
\end{definition}

\begin{theorem}[Blocked coefficients are pulled back from BNT labels]%
    \label{thm:mpdo_bnt_blocked_basis_coefficient_eq}
    \lean{MPOTensor.BNTBlockedBasisCoefficientComparison.blocked_coeff_eq}
    \leanok
    \uses{def:mpdo_bnt_blocked_basis_coefficient_comparison}
    Under a BNT comparison with blocked coefficients, each chosen blocked-basis
    multiplication coefficient is the corresponding BNT-label coefficient.
\end{theorem}
\begin{proof}\leanok
    This is exactly the defining equality of the comparison.
\end{proof}

\begin{definition}[Positive-length BNT-label trace-power form]%
    \label{def:mpdo_bnt_label_positive_length_chi_trace_power_form}
    \lean{MPOTensor.BNTLabelCoefficientFamily.HasPositiveLengthChiTracePowerForm}
    \leanok
    \uses{def:mpdo_bnt_label_coefficient_family,
        def:mpdo_diagonal_chi_family}
    A BNT-label coefficient system and a diagonal $\chi$-family are
    compatible when, for every positive length $L$,
    \[
        c^{(L)}_{\alpha,\beta,\gamma}
        = \operatorname{tr}(\chi_{\alpha,\beta,\gamma}^{L}).
    \]
    The same matrix $\chi_{\alpha,\beta,\gamma}$ is used for all positive
    $L$.
\end{definition}

\begin{theorem}[BNT-label coefficient as a positive-length trace power]%
    \label{thm:mpdo_bnt_label_chi_eq_trace_matrix_pow}
    \lean{MPOTensor.BNTLabelCoefficientFamily.HasPositiveLengthChiTracePowerForm.eq_trace_matrix_pow}
    \leanok
    \uses{def:mpdo_bnt_label_positive_length_chi_trace_power_form,
        thm:mpdo_trace_matrix_pow}
    Under positive-length BNT-label trace-power form, for every $L>0$,
    $c^{(L)}_{\alpha,\beta,\gamma}$ is the trace of
    $\chi_{\alpha,\beta,\gamma}^{L}$.
\end{theorem}
\begin{proof}\leanok
    Combine the defining positive-length trace-power identity with the trace
    formula for powers of a diagonal matrix.
\end{proof}

\begin{theorem}[Canonical BNT coefficients from a $\chi$-family]%
    \label{thm:mpdo_bnt_label_of_chi_positive_length_trace_power}
    \lean{MPOTensor.BNTLabelCoefficientFamily.ofChi_hasPositiveLengthChiTracePowerForm}
    \leanok
    \uses{def:mpdo_bnt_label_coefficient_family,
        def:mpdo_bnt_label_positive_length_chi_trace_power_form}
    The coefficient family canonically determined by a diagonal
    $\chi$-family satisfies the positive-length trace-power condition with
    respect to the same $\chi$-family.
\end{theorem}
\begin{proof}\leanok
    This is immediate from the definition of the canonical coefficient family.
\end{proof}

\begin{definition}[Positive BNT-label $\chi$ trace-power witness]%
    \label{def:mpdo_positive_bnt_label_chi_trace_power_form}
    \lean{MPOTensor.PositiveBNTLabelChiTracePowerForm}
    \lean{MPOTensor.PositiveBNTLabelChiTracePowerForm.ofChi}
    \leanok
    \uses{def:mpdo_bnt_label_positive_length_chi_trace_power_form,
        def:mpdo_diagonal_chi_family,
        thm:mpdo_bnt_label_of_chi_positive_length_trace_power}
    A positive BNT-label $\chi$ witness consists of a length-independent
    diagonal $\chi_{\alpha,\beta,\gamma}$-family, positivity of every
    diagonal entry, and the positive-length trace-power identity for the
    BNT-label coefficients. This is the coefficient statement of
    \cite[Theorem~4.14(ii)]{Cirac2016MPDO_arXiv}; constructing the witness from an MPDO tensor and
    comparing it to blocked bases remain separate obligations.
    If the coefficient family is chosen canonically from the same
    $\chi$-family, positivity of the diagonal entries alone gives such a
    witness.
\end{definition}

\begin{theorem}[Positive BNT-label witness gives trace powers]%
    \label{thm:mpdo_positive_bnt_label_chi_eq_trace_matrix_pow}
    \lean{MPOTensor.PositiveBNTLabelChiTracePowerForm.eq_trace_pow}
    \leanok
    \uses{def:mpdo_positive_bnt_label_chi_trace_power_form,
        thm:mpdo_bnt_label_chi_eq_trace_matrix_pow}
    A positive BNT-label $\chi$ witness gives, for every $L>0$, the formula
    $c^{(L)}_{\alpha,\beta,\gamma}
        =\operatorname{tr}(\chi_{\alpha,\beta,\gamma}^{L})$.
\end{theorem}
\begin{proof}\leanok
    Apply the trace reformulation of the positive-length trace-power identity
    from the witness.
\end{proof}

\begin{theorem}[Same-length BNT product with chi-trace coefficients]%
    \label{thm:mpdo_bnt_label_same_length_product_eq_sum_chi_trace_pow}
    \lean{MPOTensor.BNTLabelOperatorFamily.HasSameLengthProductForm.eq_sum_chi_trace_pow}
    \leanok
    \uses{thm:mpdo_bnt_label_same_length_product_eq_sum,
        thm:mpdo_positive_bnt_label_chi_eq_trace_matrix_pow}
    If the same-length BNT product identity holds and the BNT-label
    coefficients come from a positive length-independent $\chi$-witness, then
    for every $L>0$,
    \[
        O_L(M_\alpha)O_L(M_\beta)
        = \sum_\gamma
        \operatorname{tr}\bigl(    \chi_{\alpha,\beta,\gamma}^{\,L}
        \bigr) O_L(M_\gamma).
    \]
\end{theorem}
\begin{proof}\leanok
    Substitute the positive-length trace-power formula into the same-length BNT
    product expansion.
\end{proof}

\begin{theorem}[Same-length BNT product for canonical chi coefficients]%
    \label{thm:mpdo_bnt_label_same_length_product_eq_sum_of_chi_trace_pow}
    \lean{MPOTensor.BNTLabelOperatorFamily.HasSameLengthProductForm.eq_sum_ofChi_trace_pow}
    \leanok
    \uses{thm:mpdo_bnt_label_same_length_product_eq_sum,
        def:mpdo_bnt_label_coefficient_family}
    If the same-length BNT product identity holds with the coefficient family
    canonically determined by a diagonal $\chi$-family, then for every
    $L>0$,
    \[
        O_L(M_\alpha)O_L(M_\beta)
        = \sum_\gamma
        \operatorname{tr}\bigl(    \chi_{\alpha,\beta,\gamma}^{\,L}
        \bigr) O_L(M_\gamma).
    \]
\end{theorem}
\begin{proof}\leanok
    Substitute the canonical trace-power formula for the coefficients into the
    same-length BNT product expansion.
\end{proof}

\begin{theorem}[BNT idempotent coefficients as chi traces]%
    \label{thm:mpdo_bnt_label_idempotent_eq_sum_chi_trace}
    \lean{MPOTensor.BNTLabelTraceScalarFamily.HasIdempotentCoefficientForm.eq_sum_chi_trace}
    \leanok
    \uses{thm:mpdo_bnt_label_idempotent_coefficient_eq_sum,
        thm:mpdo_positive_bnt_label_chi_eq_trace_matrix_pow}
    If the BNT trace scalars satisfy the idempotent coefficient condition and
    the BNT-label coefficients come from a positive length-independent
    $\chi$-witness, then
    \[
        m_\gamma =
        \sum_{\alpha,\beta}
        \operatorname{tr}(\chi_{\alpha,\beta,\gamma})
        m_\alpha m_\beta.
    \]
\end{theorem}
\begin{proof}\leanok
    Substitute the length-one trace-power formula into the idempotent
    coefficient identity.
\end{proof}

\begin{theorem}[BNT idempotent coefficients for canonical chi coefficients]%
    \label{thm:mpdo_bnt_label_idempotent_eq_sum_of_chi_trace}
    \lean{MPOTensor.BNTLabelTraceScalarFamily.HasIdempotentCoefficientForm.eq_sum_ofChi_trace}
    \leanok
    \uses{thm:mpdo_bnt_label_idempotent_coefficient_eq_sum,
        def:mpdo_bnt_label_coefficient_family}
    If the BNT trace scalars satisfy the idempotent coefficient condition with
    the coefficient family canonically determined by a diagonal
    $\chi$-family, then
    \[
        m_\gamma =
        \sum_{\alpha,\beta}
        \operatorname{tr}(\chi_{\alpha,\beta,\gamma})
        m_\alpha m_\beta.
    \]
\end{theorem}
\begin{proof}\leanok
    Substitute the length-one canonical trace formula into the idempotent
    coefficient identity.
\end{proof}

\begin{theorem}[Blocked coefficients inherit BNT trace powers]%
    \label{thm:mpdo_bnt_blocked_basis_coefficient_eq_trace_pow}
    \lean{MPOTensor.BNTBlockedBasisCoefficientComparison.blocked_coeff_eq_trace_pow}
    \lean{MPOTensor.BNTBlockedBasisCoefficientComparison.%
        blocked_coeff_eq_ofChi_trace_pow}
    \lean{MPOTensor.BNTBlockedBasisCoefficientComparison.%
        blocked_coeff_eq_pulledBlockedChi_trace_pow}
    \leanok
    \uses{def:mpdo_bnt_blocked_basis_coefficient_comparison,
        def:mpdo_positive_bnt_label_chi_trace_power_form,
        thm:mpdo_positive_bnt_label_chi_eq_trace_matrix_pow}
    If the chosen blocked-basis coefficients are compared with the fixed
    BNT-label coefficient system, and that BNT-label system has a positive
    length-independent $\chi$-witness, then each positive-length blocked
    coefficient is the trace power of the corresponding BNT-label
    $\chi$-matrix:
    \[
        c^{(n)}_{i,j,k}
        = \operatorname{tr}
        \bigl(    \chi_{\sigma_n(i),\sigma_n(j),\tau_n(k)}^{\,n}
        \bigr).
    \]
\end{theorem}
\begin{proof}\leanok
    Use the comparison equality to replace the blocked-basis coefficient by the
    corresponding BNT-label coefficient, and then apply the BNT-label
    trace-power formula.
    In the special case where the compared BNT-label coefficient family is
    already the canonical family determined by $\chi$, the same conclusion
    follows directly from the canonical trace formula.
\end{proof}

\begin{definition}[Blocked $\chi$ family for multiplication coefficients]%
    \label{def:mpdo_blocked_structure_chi_family}
    \lean{MPOTensor.AlgebraStructureData.BlockedStructureChiFamily}
    \leanok
    \uses{def:mpdo_blocked_structure_coefficients}
    A dependent diagonal matrix family indexed by the blocked multiplication
    triples $(n,i,j,k)$, where $i$ and $j$ label basis elements of
    $\mathcal{A}_n$ and $k$ labels a basis element of $\mathcal{A}_{2n}$.
    This is a blocked-basis analog of the uniform BNT-label family in
    \cite[Theorem~4.14(ii)]{Cirac2016MPDO_arXiv}; the remaining uniformity
    gap is recorded in
    \path{docs/paper-gaps/cpgsv17_blocked_chi_uniformity.tex}.
\end{definition}

\begin{theorem}[Trace of a blocked $\chi$ power]%
    \label{thm:mpdo_blocked_structure_chi_trace_matrix_pow}
    \lean{MPOTensor.AlgebraStructureData.BlockedStructureChiFamily.trace_matrix_pow}
    \leanok
    \uses{def:mpdo_blocked_structure_chi_family}
    For each blocked multiplication triple $(n,i,j,k)$ and every exponent
    $L$, the trace of the $L$-th power of its diagonal matrix is the sum of the
    $L$-th powers of the corresponding diagonal entries.
\end{theorem}
\begin{proof}\leanok
    Diagonal matrices raise to powers entrywise and the trace of a diagonal
    matrix is the sum of its entries.
\end{proof}

\begin{definition}[Trace-power form for blocked multiplication coefficients]%
    \label{def:mpdo_has_blocked_structure_chi_trace_power_form}
    \lean{MPOTensor.AlgebraStructureData.HasBlockedStructureChiTracePowerForm}
    \leanok
    \uses{def:mpdo_blocked_structure_coefficients,
        def:mpdo_blocked_structure_chi_family}
    The blocked multiplication coefficients are in trace-power form when, for
    every positive blocked size $n$ and every multiplication triple $(i,j,k)$,
    the coefficient $c^{(n)}_{i,j,k}$ equals
    \[
        \sum_r \chi_{n,i,j,k,r}^{\,n}.
    \]
    Equivalently, by
    Theorem~\ref{thm:mpdo_blocked_structure_chi_trace_matrix_pow}, it is the
    trace of the $n$-th power of the corresponding diagonal matrix. The exponent
    $n$ is the source blocking length of the two factors in $\mathcal{A}_n$; the
    product is expanded in the basis of $\mathcal{A}_{2n}$.
\end{definition}

\begin{theorem}[Blocked coefficient as a trace power]%
    \label{thm:mpdo_blocked_structure_chi_eq_trace_matrix_pow}
    \lean{MPOTensor.AlgebraStructureData.HasBlockedStructureChiTracePowerForm.eq_trace_matrix_pow}
    \leanok
    \uses{def:mpdo_has_blocked_structure_chi_trace_power_form,
        thm:mpdo_blocked_structure_chi_trace_matrix_pow}
    Under trace-power form for the blocked multiplication coefficients,
    for every positive blocked size $n$, the coefficient $c^{(n)}_{i,j,k}$ is
    the trace of the $n$-th power of the diagonal matrix attached to
    $(n,i,j,k)$.
\end{theorem}
\begin{proof}\leanok
    Combine the defining trace-power identity with the trace formula for
    powers of a diagonal matrix.
\end{proof}

\begin{definition}[Positive blocked $\chi$ trace-power witness]%
    \label{def:mpdo_positive_blocked_structure_chi_trace_power_form}
    \lean{MPOTensor.AlgebraStructureData.PositiveBlockedStructureChiTracePowerForm}
    \leanok
    \uses{def:mpdo_has_blocked_structure_chi_trace_power_form,
        def:mpdo_blocked_structure_chi_family}
    A positive blocked $\chi$ trace-power witness consists of a blocked
    $\chi$ family, positivity of all its diagonal entries, and the
    positive-length trace-power identity for the blocked multiplication
    coefficients. This is the blocked-basis analog of the positive
    diagonal matrices in
    \cite[Theorem~4.14(ii)]{Cirac2016MPDO_arXiv}; the uniform BNT-label gap
    is recorded in
    \path{docs/paper-gaps/cpgsv17_blocked_chi_uniformity.tex}.
\end{definition}

\begin{theorem}[BNT witnesses give blocked $\chi$ witnesses]%
    \label{thm:mpdo_bnt_blocked_basis_to_positive_blocked_chi_trace_power_form}
    \lean{MPOTensor.BNTBlockedBasisCoefficientComparison.blockedLabel}
    \lean{MPOTensor.BNTBlockedBasisCoefficientComparison.pulledBlockedChiFamily}
    \lean{MPOTensor.BNTBlockedBasisCoefficientComparison.pulledBlockedChiFamily_toDiagonal_of_pos}
    \lean{MPOTensor.BNTBlockedBasisCoefficientComparison.pulledBlockedChi_tracePowerCoeff_of_pos}
    \lean{MPOTensor.BNTBlockedBasisCoefficientComparison.pulledBlockedChi_dim_of_pos}
    \lean{MPOTensor.BNTBlockedBasisCoefficientComparison.pulledBlockedChi_trace_matrix_pow_of_pos}
    \lean{MPOTensor.BNTBlockedBasisCoefficientComparison.%
        toPositiveBlockedStructureChiTracePowerForm}
    \leanok
    \uses{def:mpdo_bnt_blocked_basis_coefficient_comparison,
        def:mpdo_positive_bnt_label_chi_trace_power_form,
        def:mpdo_positive_blocked_structure_chi_trace_power_form,
        thm:mpdo_bnt_blocked_basis_coefficient_eq}
    Suppose that the chosen blocked-basis coefficients are compared with a
    fixed BNT-label coefficient system, and that this BNT-label system has a
    positive length-independent $\chi$-witness. Then the chosen blocked
    bases carry a positive blocked $\chi$ trace-power witness.
    The construction is obtained by pulling back the BNT-label
    $\chi_{\alpha,\beta,\gamma}$-matrices along the source and target label
    maps.
    % Formalization note: the zero-length component of the blocked chi family is
    % filled by empty diagonal matrices; it plays no role in the positive-length
    % trace-power statement.
\end{theorem}
\begin{proof}\leanok
    Reindex the positive BNT-label $\chi$-family along the blocked-basis
    label map. Positivity is preserved by reindexing, and the trace-power
    identity follows by substituting the blocked-basis comparison into the
    BNT-label trace-power formula.
\end{proof}

\begin{definition}[BNT-label theorem data]%
    \label{def:mpdo_bnt_label_theorem_data}
    \lean{MPOTensor.BNTLabelTheoremData}
    \lean{MPOTensor.BNTLabelTheoremData.same_length_product_form}
    \lean{MPOTensor.BNTLabelTheoremData.idempotent_coefficient_form}
    \lean{MPOTensor.BNTLabelTheoremData.same_length_product_form_ofChi}
    \lean{MPOTensor.BNTLabelTheoremData.idempotent_coefficient_form_ofChi}
    \lean{MPOTensor.BNTLabelTheoremData.same_length_product_eq_sum_ofChi}
    \lean{MPOTensor.BNTLabelTheoremData.idempotent_eq_sum_ofChi}
    \lean{MPOTensor.BNTLabelTheoremData.positive_chi_pos_entries}
    \lean{MPOTensor.BNTLabelTheoremData.positive_chi_trace_power}
    \lean{MPOTensor.BNTLabelTheoremData.labelAssignment}
    \lean{MPOTensor.BNTLabelTheoremData.sourceLabel}
    \lean{MPOTensor.BNTLabelTheoremData.targetLabel}
    \lean{MPOTensor.BNTLabelTheoremData.positiveBlockedChi}
    \lean{MPOTensor.BNTLabelTheoremData.ofChi}
    \leanok
    \uses{def:mpdo_bnt_label_coefficient_family,
        def:mpdo_bnt_label_operator_family,
        def:mpdo_bnt_label_trace_scalar_family,
        def:mpdo_bnt_label_same_length_product_form,
        def:mpdo_bnt_label_idempotent_coefficient_form,
        def:mpdo_positive_bnt_label_chi_trace_power_form,
        def:mpdo_bnt_blocked_basis_label_assignment,
        def:mpdo_bnt_blocked_basis_coefficient_comparison}
    The BNT-label theorem data consist of the fixed BNT-label coefficient
    system, the same-length BNT operator family and product identity, the trace
    scalars and idempotent condition, the positive length-independent
    $\chi$-witness, and the comparison with the chosen blocked bases.
    They also determine the pulled-back blocked-basis $\chi$-family. These
    data are the source-side hypotheses from which the blocked-basis
    consequences below are derived.
    In the source-side special case where the coefficient family is
    canonically determined by the same $\chi$-family, theorem data are built
    from that $\chi$-family, its positivity, and the product, idempotent, and
    blocked-comparison statements for the resulting canonical coefficients.
    The same product and idempotent predicates may also be rephrased using the
    canonical coefficient family determined by the $\chi$-matrices carried by
    the data, with the corresponding displayed equations.
\end{definition}

\begin{theorem}[BNT theorem data give the same-length product law]%
    \label{thm:mpdo_bnt_label_theorem_data_same_length_product_eq_sum}
    \lean{MPOTensor.BNTLabelTheoremData.same_length_product_eq_sum}
    \leanok
    \uses{def:mpdo_bnt_label_theorem_data,
        thm:mpdo_bnt_label_same_length_product_eq_sum}
    BNT-label theorem data give the same-length product equation
    \[
        O_L(M_\alpha)O_L(M_\beta)
        = \sum_\gamma
        c^{(L)}_{\alpha,\beta,\gamma} O_L(M_\gamma)
    \]
    for every positive length $L$.
\end{theorem}
\begin{proof}\leanok
    This is the same-length product law carried by the theorem data.
\end{proof}

\begin{theorem}[BNT theorem data give the idempotent scalar law]%
    \label{thm:mpdo_bnt_label_theorem_data_idempotent_eq_sum}
    \lean{MPOTensor.BNTLabelTheoremData.idempotent_eq_sum}
    \leanok
    \uses{def:mpdo_bnt_label_theorem_data,
        thm:mpdo_bnt_label_idempotent_coefficient_eq_sum}
    BNT-label theorem data give the idempotent scalar equation
    \[
        m_\gamma =
        \sum_{\alpha,\beta}
        c^{(1)}_{\alpha,\beta,\gamma}m_\alpha m_\beta.
    \]
\end{theorem}
\begin{proof}\leanok
    This is the idempotent coefficient law carried by the theorem data.
\end{proof}

\begin{theorem}[BNT theorem data give the blocked coefficient comparison]%
    \label{thm:mpdo_bnt_label_theorem_data_blocked_coeff_eq}
    \lean{MPOTensor.BNTLabelTheoremData.blocked_coeff_eq}
    \lean{MPOTensor.BNTLabelTheoremData.blockedComparison_ofChi}
    \lean{MPOTensor.BNTLabelTheoremData.blocked_coeff_eq_ofChi}
    \leanok
    \uses{def:mpdo_bnt_label_theorem_data,
        thm:mpdo_bnt_blocked_basis_coefficient_eq,
        thm:mpdo_bnt_label_theorem_data_coeff_eq_trace_pow}
    BNT-label theorem data give the comparison between the blocked-basis
    coefficients and the BNT-label coefficients through the source and target
    label maps. Since the positive-length coefficients agree with the
    canonical coefficient family determined by the same $\chi$-matrices, the
    comparison can also be written with those canonical coefficients.
\end{theorem}
\begin{proof}\leanok
    This is the blocked-basis coefficient comparison carried by the theorem
    data.
\end{proof}

\begin{theorem}[BNT theorem data give coefficient trace powers]%
    \label{thm:mpdo_bnt_label_theorem_data_coeff_eq_trace_pow}
    \lean{MPOTensor.BNTLabelTheoremData.coeff_eq_trace_pow}
    \lean{MPOTensor.BNTLabelTheoremData.coeff_eq_ofChi_coeff}
    \leanok
    \uses{def:mpdo_bnt_label_theorem_data,
        thm:mpdo_positive_bnt_label_chi_eq_trace_matrix_pow}
    BNT-label theorem data give the coefficient identity
    $c^{(L)}_{\alpha,\beta,\gamma}
        = \operatorname{tr}(\chi_{\alpha,\beta,\gamma}^{L})$
    for every positive length $L$. Equivalently, at every positive length, the
    coefficient family agrees with the canonical coefficient family determined
    by the same $\chi$-matrices.
\end{theorem}
\begin{proof}\leanok
    This is the positive-length $\chi$-trace identity belonging to the
    positive BNT-label $\chi$-witness in the theorem data.
\end{proof}

\begin{theorem}[BNT theorem data give positive chi entries]%
    \label{thm:mpdo_bnt_label_theorem_data_chi_entry_pos}
    \lean{MPOTensor.BNTLabelTheoremData.chi_entry_pos}
    \leanok
    \uses{def:mpdo_bnt_label_theorem_data,
        def:mpdo_positive_bnt_label_chi_trace_power_form}
    BNT-label theorem data give positivity of every diagonal entry of
    $\chi_{\alpha,\beta,\gamma}$.
\end{theorem}
\begin{proof}\leanok
    This is the positivity condition carried by the positive BNT-label
    $\chi$-witness in the theorem data.
\end{proof}

\begin{theorem}[BNT theorem data give the chi-trace product law]%
    \label{thm:mpdo_bnt_label_theorem_data_same_length_product_eq_sum_chi_trace_pow}
    \lean{MPOTensor.BNTLabelTheoremData.same_length_product_eq_sum_chi_trace_pow}
    \leanok
    \uses{def:mpdo_bnt_label_theorem_data,
        thm:mpdo_bnt_label_same_length_product_eq_sum_chi_trace_pow}
    BNT-label theorem data give the same-length product expansion with
    coefficients written as traces of powers of the length-independent
    $\chi_{\alpha,\beta,\gamma}$-matrices.
\end{theorem}
\begin{proof}\leanok
    Apply the chi-trace product law to the product identity and the positive
    BNT-label $\chi$-witness belonging to the theorem data.
\end{proof}

\begin{theorem}[BNT theorem data give the chi-trace idempotent law]%
    \label{thm:mpdo_bnt_label_theorem_data_idempotent_eq_sum_chi_trace}
    \lean{MPOTensor.BNTLabelTheoremData.idempotent_eq_sum_chi_trace}
    \leanok
    \uses{def:mpdo_bnt_label_theorem_data,
        thm:mpdo_bnt_label_idempotent_eq_sum_chi_trace}
    BNT-label theorem data give the idempotent scalar identity with the
    length-one coefficients written as traces of the corresponding
    $\chi$-matrices.
\end{theorem}
\begin{proof}\leanok
    Apply the chi-trace idempotent law to the idempotent condition and the
    positive BNT-label $\chi$-witness belonging to the theorem data.
\end{proof}

\begin{theorem}[BNT theorem data give blocked coefficient trace powers]%
    \label{thm:mpdo_bnt_label_theorem_data_blocked_coeff_eq_trace_pow}
    \lean{MPOTensor.BNTLabelTheoremData.blocked_coeff_eq_trace_pow}
    \leanok
    \uses{def:mpdo_bnt_label_theorem_data,
        thm:mpdo_bnt_blocked_basis_coefficient_eq_trace_pow}
    BNT-label theorem data give the blocked-basis coefficient trace-power
    formula obtained by pulling the BNT-label
    $\chi_{\alpha,\beta,\gamma}$-matrices back along the blocked-basis
    comparison maps.
\end{theorem}
\begin{proof}\leanok
    Apply the blocked-basis coefficient trace-power theorem to the comparison
    and the positive BNT-label $\chi$-witness belonging to the theorem data.
\end{proof}

\begin{theorem}[BNT theorem data give a positive blocked $\chi$ witness]%
    \label{thm:mpdo_bnt_label_theorem_data_to_positive_blocked_chi_trace_power_form}
    \lean{MPOTensor.BNTLabelTheoremData.toPositiveBlockedStructureChiTracePowerForm}
    \lean{MPOTensor.BNTLabelTheoremData.positiveBlockedChi_toDiagonal_of_pos}
    \lean{MPOTensor.BNTLabelTheoremData.positiveBlockedChi_tracePowerCoeff_of_pos}
    \lean{MPOTensor.BNTLabelTheoremData.positiveBlockedChi_dim_of_pos}
    \lean{MPOTensor.BNTLabelTheoremData.positiveBlockedChi_trace_matrix_pow_of_pos}
    \leanok
    \uses{def:mpdo_bnt_label_theorem_data,
        thm:mpdo_bnt_blocked_basis_to_positive_blocked_chi_trace_power_form}
    BNT-label theorem data give a positive blocked-basis $\chi$ trace-power
    witness by pulling the length-independent BNT-label $\chi$-family back
    along the blocked-basis comparison maps.
\end{theorem}
\begin{proof}\leanok
    Apply the BNT-to-blocked witness construction to the comparison and the
    positive BNT-label $\chi$-witness belonging to the theorem data.
\end{proof}

\begin{theorem}[BNT theorem data give blocked trace-power form]%
    \label{thm:mpdo_bnt_label_theorem_data_positive_blocked_chi_trace_power}
    \lean{MPOTensor.BNTLabelTheoremData.positiveBlockedChi_tracePower}
    \leanok
    \uses{def:mpdo_bnt_label_theorem_data,
        thm:mpdo_bnt_label_theorem_data_to_positive_blocked_chi_trace_power_form}
    The pulled-back blocked-basis $\chi$-family obtained from BNT-label
    theorem data satisfies the blocked trace-power predicate.
\end{theorem}
\begin{proof}\leanok
    This is the trace-power part of the positive blocked-basis
    $\chi$-witness obtained from the theorem data.
\end{proof}

\begin{theorem}[BNT theorem data give positive blocked chi families]%
    \label{thm:mpdo_bnt_label_theorem_data_positive_blocked_chi_pos_entries}
    \lean{MPOTensor.BNTLabelTheoremData.positiveBlockedChi_posEntries}
    \leanok
    \uses{def:mpdo_bnt_label_theorem_data,
        thm:mpdo_bnt_label_theorem_data_to_positive_blocked_chi_trace_power_form}
    The pulled-back blocked-basis $\chi$-family obtained from BNT-label
    theorem data has positive diagonal entries.
\end{theorem}
\begin{proof}\leanok
    This is the positivity part of the positive blocked-basis
    $\chi$-witness obtained from the theorem data.
\end{proof}

\begin{theorem}[BNT theorem data give positive pulled-back blocked entries]%
    \label{thm:mpdo_bnt_label_theorem_data_positive_blocked_chi_entry_pos}
    \lean{MPOTensor.BNTLabelTheoremData.positiveBlockedChi_entry_pos}
    \leanok
    \uses{def:mpdo_bnt_label_theorem_data,
        thm:mpdo_bnt_label_theorem_data_to_positive_blocked_chi_trace_power_form}
    The pulled-back blocked-basis $\chi$-entries obtained from BNT-label
    theorem data are positive.
\end{theorem}
\begin{proof}\leanok
    This is the positivity part of the positive blocked-basis
    $\chi$-witness obtained from the theorem data.
\end{proof}

\begin{theorem}[BNT theorem data give pulled-back blocked trace powers]%
    \label{thm:mpdo_bnt_label_theorem_data_blocked_coeff_eq_positive_blocked_chi_trace_pow}
    \lean{MPOTensor.BNTLabelTheoremData.blocked_coeff_eq_positiveBlockedChi_trace_pow}
    \leanok
    \uses{def:mpdo_bnt_label_theorem_data,
        thm:mpdo_bnt_label_theorem_data_to_positive_blocked_chi_trace_power_form,
        thm:mpdo_positive_blocked_structure_chi_eq_trace_matrix_pow}
    The positive-length blocked-basis coefficients are traces of powers of the
    pulled-back blocked-basis $\chi$-matrices obtained from BNT-label
    theorem data.
\end{theorem}
\begin{proof}\leanok
    Apply the trace-power identity of the positive blocked-basis
    $\chi$-witness obtained from the theorem data.
\end{proof}

\begin{definition}[Existential BNT-label theorem witness]%
    \label{def:mpdo_bnt_label_theorem_witness}
    \lean{MPOTensor.BNTLabelTheoremWitness}
    \lean{MPOTensor.HasBNTLabelTheoremWitness}
    \lean{MPOTensor.BNTLabelTheoremWitness.toTheoremData}
    \lean{MPOTensor.BNTLabelTheoremWitness.coeffs}
    \lean{MPOTensor.BNTLabelTheoremWitness.operators}
    \lean{MPOTensor.BNTLabelTheoremWitness.traceScalars}
    \lean{MPOTensor.BNTLabelTheoremWitness.positiveChi}
    \lean{MPOTensor.BNTLabelTheoremWitness.blockedComparison}
    \lean{MPOTensor.BNTLabelTheoremWitness.same_length_product_form}
    \lean{MPOTensor.BNTLabelTheoremWitness.idempotent_coefficient_form}
    \lean{MPOTensor.BNTLabelTheoremWitness.same_length_product_form_ofChi}
    \lean{MPOTensor.BNTLabelTheoremWitness.idempotent_coefficient_form_ofChi}
    \lean{MPOTensor.BNTLabelTheoremWitness.same_length_product_eq_sum_ofChi}
    \lean{MPOTensor.BNTLabelTheoremWitness.idempotent_eq_sum_ofChi}
    \lean{MPOTensor.BNTLabelTheoremWitness.positive_chi_pos_entries}
    \lean{MPOTensor.BNTLabelTheoremWitness.positive_chi_trace_power}
    \lean{MPOTensor.BNTLabelTheoremWitness.labelAssignment}
    \lean{MPOTensor.BNTLabelTheoremWitness.sourceLabel}
    \lean{MPOTensor.BNTLabelTheoremWitness.targetLabel}
    \lean{MPOTensor.BNTLabelTheoremWitness.positiveBlockedChi}
    \lean{MPOTensor.BNTLabelTheoremWitness.ofChi}
    \lean{MPOTensor.BNTLabelTheoremWitness.hasBNTLabelTheoremWitness}
    \lean{MPOTensor.BNTLabelTheoremWitness.ofChi_hasBNTLabelTheoremWitness}
    \leanok
    \uses{def:mpdo_bnt_label_theorem_data}
    An existential BNT-label theorem witness consists of the finite BNT-label
    type, the same-length operator spaces, their algebraic structure, and the
    corresponding BNT-label theorem data, including the pulled-back
    blocked-basis $\chi$-family. The construction of such a witness from an
    MPDO tensor is the outstanding step toward \cite[Theorem~4.14(ii)]{Cirac2016MPDO_arXiv}.
    The proposition-level form is the nonemptiness of this witness type.
    In the source-side special case where the coefficients are canonically
    determined by a $\chi$-family, the witness is obtained from
    that same $\chi$-family together with the remaining product, idempotent,
    and blocked-comparison statements.
    For any such witness, the product and idempotent predicates can be
    rephrased with the canonical coefficient family determined by its
    $\chi$-matrices, with the corresponding displayed equations.
\end{definition}

\begin{theorem}[BNT theorem witnesses give source equations]%
    \label{thm:mpdo_has_bnt_label_theorem_witness_source_equations}
    \lean{MPOTensor.HasBNTLabelTheoremWitness.exists_source_predicates}
    \lean{MPOTensor.HasBNTLabelTheoremWitness.exists_positive_length_coeff_eq_ofChi}
    \lean{MPOTensor.HasBNTLabelTheoremWitness.exists_source_ofChi_predicates}
    \lean{MPOTensor.HasBNTLabelTheoremWitness.exists_source_ofChi_equations}
    \lean{MPOTensor.HasBNTLabelTheoremWitness.exists_source_coefficient_equations}
    \lean{MPOTensor.HasBNTLabelTheoremWitness.exists_source_chi_trace_equations}
    \leanok
    \uses{def:mpdo_bnt_label_theorem_witness}
    A proposition-level BNT-label theorem witness contains the source-side
    same-length product law, the idempotent scalar law, the positive-length
    trace-power law, and positivity of the diagonal
    $\chi_{\alpha,\beta,\gamma}$-entries. After unpacking the witness, the
    same-length product and idempotent equations can first be written with the
    BNT-label coefficients $c^{(L)}_{\alpha,\beta,\gamma}$, and then with
    the coefficients expressed as
    $\operatorname{tr}(\chi_{\alpha,\beta,\gamma}^{\,L})$ and
    $\operatorname{tr}(\chi_{\alpha,\beta,\gamma})$, respectively. At
    every positive length, the BNT-label coefficient family also agrees with
    the canonical coefficient family determined by these $\chi$-matrices.
    Consequently, the product and idempotent predicates may be stated directly
    using that canonical coefficient family, and the same displayed source
    equations hold with those canonical coefficients.
\end{theorem}
\begin{proof}\leanok
    Unpack the existential witness and apply the corresponding witness-level
    equations.
\end{proof}

\begin{theorem}[BNT theorem witnesses give coefficient trace powers]%
    \label{thm:mpdo_bnt_label_theorem_witness_coeff_eq_trace_pow}
    \lean{MPOTensor.BNTLabelTheoremWitness.coeff_eq_trace_pow}
    \lean{MPOTensor.BNTLabelTheoremWitness.coeff_eq_ofChi_coeff}
    \leanok
    \uses{def:mpdo_bnt_label_theorem_witness,
        thm:mpdo_bnt_label_theorem_data_coeff_eq_trace_pow}
    An existential BNT-label theorem witness gives the formula
    $c_{\alpha,\beta,\gamma}^{(L)}
        = \operatorname{tr}(\chi_{\alpha,\beta,\gamma}^{\,L})$
    for every positive length $L$. Equivalently, at positive lengths, its
    coefficient family agrees with the canonical coefficient family determined
    by the same $\chi$-matrices.
\end{theorem}
\begin{proof}\leanok
    Apply the corresponding theorem-data identity carried by the witness.
\end{proof}

\begin{theorem}[BNT theorem witnesses give positive chi entries]%
    \label{thm:mpdo_bnt_label_theorem_witness_chi_entry_pos}
    \lean{MPOTensor.BNTLabelTheoremWitness.chi_entry_pos}
    \leanok
    \uses{def:mpdo_bnt_label_theorem_witness,
        def:mpdo_positive_bnt_label_chi_trace_power_form}
    An existential BNT-label theorem witness gives positivity of every
    diagonal entry of $\chi_{\alpha,\beta,\gamma}$.
\end{theorem}
\begin{proof}\leanok
    This follows from the positivity condition in the positive BNT-label
    $\chi$-witness.
\end{proof}

\begin{theorem}[BNT theorem witnesses give the same-length product law]%
    \label{thm:mpdo_bnt_label_theorem_witness_same_length_product_eq_sum}
    \lean{MPOTensor.BNTLabelTheoremWitness.same_length_product_eq_sum}
    \leanok
    \uses{def:mpdo_bnt_label_theorem_witness,
        thm:mpdo_bnt_label_theorem_data_same_length_product_eq_sum}
    An existential BNT-label theorem witness gives the product identity
    $O_L(M_\alpha)O_L(M_\beta)
        = \sum_\gamma c_{\alpha,\beta,\gamma}^{(L)} O_L(M_\gamma)$
    for every positive length $L$.
\end{theorem}
\begin{proof}\leanok
    Apply the corresponding theorem-data identity carried by the witness.
\end{proof}

\begin{theorem}[BNT theorem witnesses give the idempotent scalar law]%
    \label{thm:mpdo_bnt_label_theorem_witness_idempotent_eq_sum}
    \lean{MPOTensor.BNTLabelTheoremWitness.idempotent_eq_sum}
    \leanok
    \uses{def:mpdo_bnt_label_theorem_witness,
        thm:mpdo_bnt_label_theorem_data_idempotent_eq_sum}
    An existential BNT-label theorem witness gives the idempotent scalar
    identity
    $m_\gamma =
        \sum_{\alpha,\beta} c_{\alpha,\beta,\gamma}^{(1)}m_\alpha m_\beta$.
\end{theorem}
\begin{proof}\leanok
    Apply the corresponding theorem-data identity carried by the witness.
\end{proof}

\begin{theorem}[BNT theorem witnesses give the blocked coefficient comparison]%
    \label{thm:mpdo_bnt_label_theorem_witness_blocked_coeff_eq}
    \lean{MPOTensor.BNTLabelTheoremWitness.blocked_coeff_eq}
    \lean{MPOTensor.BNTLabelTheoremWitness.blockedComparison_ofChi}
    \lean{MPOTensor.BNTLabelTheoremWitness.blocked_coeff_eq_ofChi}
    \lean{MPOTensor.HasBNTLabelTheoremWitness.exists_blocked_coefficient_comparison}
    \lean{MPOTensor.HasBNTLabelTheoremWitness.%
        exists_blocked_coefficient_comparison_ofChi}
    \leanok
    \uses{def:mpdo_bnt_label_theorem_witness,
        thm:mpdo_bnt_label_theorem_data_blocked_coeff_eq}
    An existential BNT-label theorem witness gives the comparison between the
    blocked-basis coefficients and the BNT-label coefficients through the
    source and target label maps. The same comparison is also available with
    the canonical coefficient family determined by the witness's
    $\chi$-matrices. The proposition-level nonempty form of the witness
    gives the same comparisons after choosing a witness.
\end{theorem}
\begin{proof}\leanok
    Apply the corresponding theorem-data comparison carried by the witness.
\end{proof}

\begin{theorem}[BNT theorem witnesses give the chi-trace product law]%
    \label{thm:mpdo_bnt_label_theorem_witness_same_length_product_eq_sum_chi_trace_pow}
    \lean{MPOTensor.BNTLabelTheoremWitness.same_length_product_eq_sum_chi_trace_pow}
    \leanok
    \uses{def:mpdo_bnt_label_theorem_witness,
        thm:mpdo_bnt_label_theorem_data_same_length_product_eq_sum_chi_trace_pow}
    An existential BNT-label theorem witness gives, for every positive length
    $L$, the identity
    $O_L(M_\alpha)O_L(M_\beta)
        = \sum_\gamma
        \operatorname{tr}(\chi_{\alpha,\beta,\gamma}^{\,L}) O_L(M_\gamma)$.
\end{theorem}
\begin{proof}\leanok
    Apply the corresponding theorem-data identity carried by the witness.
\end{proof}

\begin{theorem}[BNT theorem witnesses give the chi-trace idempotent law]%
    \label{thm:mpdo_bnt_label_theorem_witness_idempotent_eq_sum_chi_trace}
    \lean{MPOTensor.BNTLabelTheoremWitness.idempotent_eq_sum_chi_trace}
    \leanok
    \uses{def:mpdo_bnt_label_theorem_witness,
        thm:mpdo_bnt_label_theorem_data_idempotent_eq_sum_chi_trace}
    An existential BNT-label theorem witness gives the idempotent scalar
    identity
    $m_\gamma =
        \sum_{\alpha,\beta}
        \operatorname{tr}(\chi_{\alpha,\beta,\gamma})m_\alpha m_\beta$.
\end{theorem}
\begin{proof}\leanok
    Apply the corresponding theorem-data identity carried by the witness.
\end{proof}

\begin{theorem}[BNT theorem witnesses give blocked coefficient trace powers]%
    \label{thm:mpdo_bnt_label_theorem_witness_blocked_coeff_eq_trace_pow}
    \lean{MPOTensor.BNTLabelTheoremWitness.blocked_coeff_eq_trace_pow}
    \leanok
    \uses{def:mpdo_bnt_label_theorem_witness,
        thm:mpdo_bnt_label_theorem_data_blocked_coeff_eq_trace_pow}
    An existential BNT-label theorem witness gives the blocked-basis
    coefficient trace-power formula obtained by pulling the BNT-label
    $\chi$-matrices back along the blocked-basis comparison maps.
\end{theorem}
\begin{proof}\leanok
    Apply the corresponding theorem-data identity carried by the witness.
\end{proof}

\begin{theorem}[BNT theorem witnesses give positive blocked $\chi$ witnesses]%
    \label{thm:mpdo_bnt_label_theorem_witness_to_positive_blocked_chi_trace_power_form}
    \lean{MPOTensor.BNTLabelTheoremWitness.toPositiveBlockedStructureChiTracePowerForm}
    \lean{MPOTensor.BNTLabelTheoremWitness.positiveBlockedChi_toDiagonal_of_pos}
    \lean{MPOTensor.BNTLabelTheoremWitness.positiveBlockedChi_tracePowerCoeff_of_pos}
    \lean{MPOTensor.BNTLabelTheoremWitness.positiveBlockedChi_dim_of_pos}
    \lean{MPOTensor.BNTLabelTheoremWitness.positiveBlockedChi_trace_matrix_pow_of_pos}
    \lean{MPOTensor.HasBNTLabelTheoremWitness.positive_blocked_chi_witness}
    \lean{MPOTensor.HasBNTLabelTheoremWitness.exists_blocked_chi_trace_power_form}
    \lean{MPOTensor.HasBNTLabelTheoremWitness.exists_blocked_chi_pullback}
    \lean{MPOTensor.HasBNTLabelTheoremWitness.exists_blocked_coeff_eq_trace_pow}
    \leanok
    \uses{def:mpdo_bnt_label_theorem_witness,
        thm:mpdo_bnt_label_theorem_data_to_positive_blocked_chi_trace_power_form}
    An existential BNT-label theorem witness gives a positive blocked-basis
    $\chi$ trace-power witness for the chosen algebra-structure data. The
    blocked-basis $\chi$-family is the pullback of the BNT-label
    $\chi$-family along the source and target comparison maps. The
    proposition-level nonempty form of the witness gives the corresponding
    existential blocked-basis trace-power family, both in coefficient form and
    in matrix-trace form.
\end{theorem}
\begin{proof}\leanok
    Apply the corresponding theorem-data consequence carried by the witness.
\end{proof}

\begin{theorem}[BNT theorem witnesses give blocked trace-power form]%
    \label{thm:mpdo_bnt_label_theorem_witness_positive_blocked_chi_trace_power}
    \lean{MPOTensor.BNTLabelTheoremWitness.positiveBlockedChi_tracePower}
    \leanok
    \uses{def:mpdo_bnt_label_theorem_witness,
        thm:mpdo_bnt_label_theorem_witness_to_positive_blocked_chi_trace_power_form}
    The pulled-back blocked-basis $\chi$-family obtained from an existential
    BNT-label theorem witness satisfies the blocked trace-power predicate.
\end{theorem}
\begin{proof}\leanok
    This is the trace-power part of the positive blocked-basis
    $\chi$-witness obtained from the existential witness.
\end{proof}

\begin{theorem}[BNT theorem witnesses give positive blocked chi families]%
    \label{thm:mpdo_bnt_label_theorem_witness_positive_blocked_chi_pos_entries}
    \lean{MPOTensor.BNTLabelTheoremWitness.positiveBlockedChi_posEntries}
    \leanok
    \uses{def:mpdo_bnt_label_theorem_witness,
        thm:mpdo_bnt_label_theorem_witness_to_positive_blocked_chi_trace_power_form}
    The pulled-back blocked-basis $\chi$-family obtained from an existential
    BNT-label theorem witness has positive diagonal entries.
\end{theorem}
\begin{proof}\leanok
    This is the positivity part of the positive blocked-basis
    $\chi$-witness obtained from the existential witness.
\end{proof}

\begin{theorem}[BNT theorem witnesses give positive pulled-back blocked entries]%
    \label{thm:mpdo_bnt_label_theorem_witness_positive_blocked_chi_entry_pos}
    \lean{MPOTensor.BNTLabelTheoremWitness.positiveBlockedChi_entry_pos}
    \leanok
    \uses{def:mpdo_bnt_label_theorem_witness,
        thm:mpdo_bnt_label_theorem_witness_to_positive_blocked_chi_trace_power_form}
    The pulled-back blocked-basis $\chi$-entries obtained from an
    existential BNT-label theorem witness are positive.
\end{theorem}
\begin{proof}\leanok
    This is the positivity part of the positive blocked-basis
    $\chi$-witness obtained from the existential witness.
\end{proof}

\begin{theorem}[BNT theorem witnesses give pulled-back blocked trace powers]%
    \label{thm:mpdo_bnt_label_theorem_witness_blocked_coeff_eq_positive_blocked_chi_trace_pow}
    \lean{MPOTensor.BNTLabelTheoremWitness.blocked_coeff_eq_positiveBlockedChi_trace_pow}
    \leanok
    \uses{def:mpdo_bnt_label_theorem_witness,
        thm:mpdo_bnt_label_theorem_witness_to_positive_blocked_chi_trace_power_form,
        thm:mpdo_positive_blocked_structure_chi_eq_trace_matrix_pow}
    The positive-length blocked-basis coefficients are traces of powers of the
    pulled-back blocked-basis $\chi$-matrices obtained from an existential
    BNT-label theorem witness.
\end{theorem}
\begin{proof}\leanok
    Apply the trace-power identity of the positive blocked-basis
    $\chi$-witness obtained from the existential witness.
\end{proof}

\begin{theorem}[Positive blocked $\chi$ witness gives trace powers]%
    \label{thm:mpdo_positive_blocked_structure_chi_eq_trace_matrix_pow}
    \lean{MPOTensor.AlgebraStructureData.PositiveBlockedStructureChiTracePowerForm.eq_trace_pow}
    \leanok
    \uses{def:mpdo_positive_blocked_structure_chi_trace_power_form,
        thm:mpdo_blocked_structure_chi_eq_trace_matrix_pow}
    A positive blocked $\chi$ trace-power witness gives, for every positive
    blocked size $n$, the formula
    $c^{(n)}_{i,j,k}=\operatorname{tr}(\chi_{n,i,j,k}^{\,n})$.
\end{theorem}
\begin{proof}\leanok
    Apply the trace reformulation of the trace-power identity from the witness.
\end{proof}
